\section{Motivation}

Model multiplicity is especially relevant for cyber defense because intrusion-detection models are not used as isolated benchmark artifacts. Their outputs may trigger alerts, change incident priorities, guide analyst attention, or become part of automated response pipelines. In these settings, two models with nearly identical aggregate performance can still imply different operational actions for the same network event. This makes predictive disagreement a security-relevant property rather than a purely statistical curiosity.

\subsection{Beyond Aggregate Detection Performance}

Intrusion-detection research commonly reports accuracy, precision, recall, F1 score, or related aggregate metrics. These metrics are indispensable, but they compress model behavior over a test set and can hide pointwise disagreement. The BA results motivating this work showed that similarly accurate neural networks can disagree on individual samples and that this disagreement can be measured through ambiguity and conflict ratios. In a cyber defense workflow, such points are precisely the cases where model choice matters: a flow classified as benign by one high-performing model and malicious by another may determine whether an analyst investigates an alert or whether an automated defense action is taken.

This concern is amplified by the cost structure of intrusion detection. False negatives may leave attacks undetected, while false positives can overwhelm analysts or desensitize operators to alerts. If a Rashomon set contains models with similar test performance but different local decisions, selecting a single model by a marginal metric advantage can make local security outcomes appear more determined than they actually are. Studying multiplicity therefore exposes where the learned decision boundary is stable and where the data and training procedure leave room for materially different decisions.

\subsection{Explanation Reliability for Security Analysts}

Explainability is frequently proposed as a way to make ML-based intrusion detection more trustworthy and usable for analysts. Prior cyber-security XAI work argues that black-box ML can reduce user confidence and that explanations are needed to help humans understand, trust, and manage ML-based defenses~\cite{zhang2022explainable}. Work on explainable intrusion detection similarly treats explanations as a mechanism for interpreting and validating alerts~\cite{szczepanski2020achieving}. However, explanations are usually generated for one selected model. The BA study showed that this can be misleading: even when accuracy is identical, SHAP-based explanations may differ substantially, and for real-world datasets predictive conflict can correlate with SHAP-value variance for important features.

This matters in cyber defense because explanations are not merely descriptive. They can influence whether analysts regard an alert as credible, which traffic characteristics they inspect, and which mitigation path they choose. If one model explains an attack prediction through failed-login features while another emphasizes service-level error rates or host-level traffic patterns, the resulting investigation may follow different hypotheses. A single-model explanation can therefore overstate the stability of the evidence behind an alert. Measuring explanation variability across similarly accurate models provides a more honest view of which explanatory signals are robust and which are artifacts of model selection.

\subsection{Adversarial and Operational Robustness}

Cyber defense also differs from many standard classification settings because the input distribution is adversarial and non-stationary. Attackers adapt, benign network behavior changes, and deployment environments rarely match benchmark assumptions. Prior work on adversarial machine learning for network intrusion detection emphasizes that ML-based NIDS can themselves become attack targets and that robustness, explainability, and security must be considered alongside predictive performance~\cite{pawlicki2020defending,choras2020machine}. This creates a natural motivation for analyzing not only whether a model performs well on average, but also whether its decisions and explanations remain stable under plausible model variation.

Recent work has begun to frame model multiplicity itself as useful for defense: Behfar and Mortier propose maintaining multiple independently evolving models and using divergence between them as a signal for adversarial behavior in distributed edge training~\cite{behfar2026model}. Although their setting differs from network intrusion detection, the underlying idea is relevant: disagreement among plausible models can carry operational information. In intrusion detection, multiplicity analysis can identify samples, attack categories, or features where the detector is underdetermined, brittle, or especially sensitive to training variation.

Model stealing provides a complementary adversarial perspective. Baltz, Hara, and Aivodji argue that high-fidelity surrogate models obtained through query access should not automatically be treated as deployment-equivalent copies of the target model~\cite{baltz2026model}. Because the extraction process only partially supervises the target's input-output behavior, many near-optimal surrogates may agree with the target at similar rates while differing in ambiguity, discrepancy, Rashomon capacity, and fairness-relevant behavior. For cyber defense, this is important in both directions: an attacker may not obtain an operationally equivalent detector even when label agreement is high, and a defender should not assume that a single fidelity or aggregate-performance measure captures the security consequences of model replacement, replication, or explanation.

\subsection{Research Gap}

Existing work has studied predictive multiplicity and explanation multiplicity in general ML settings~\cite{black2022model,marx2020predictive,brunet2022implications,laberge2023partial}, and recent model-stealing work shows that multiplicity can change how adversarial replication should be interpreted~\cite{baltz2026model}. Cyber-security research has also studied explainable intrusion detection and adversarial robustness. The intersection remains underexplored. In particular, there is limited evidence on how Rashomon-set disagreement and explanation variability interact for intrusion-detection data.

This paper addresses that gap by transferring the BA-style multiplicity and SHAP-variability analysis to NSL-KDD. The goal is not only to report whether several neural intrusion detectors achieve similar F1 scores, but to identify where they disagree, whether those disagreements coincide with unstable explanations, and what that implies for trustworthy cyber defense workflows.
